\section{MDP 基础回顾}

本讲由助教 Anakiette 主讲，作为课程总复习，从 MDP 基础出发回顾 Q-learning 的完整理论链条。

\textit{``We're going to start off by just doing a very very brief overview of MDPs.''} 这句开场既是提示，也是学习节奏的指南，提醒我们在完整的理论链条中先建立形式化符号体系，再逐层深入细节。

\subsection{马尔可夫决策过程}

\begin{importantbox}{MDP 五元组}
MDP 由 $(\mathcal{S}, \mathcal{A}, P, R, \gamma)$ 定义：
\begin{itemize}
    \item $\mathcal{S}$：状态空间
    \item $\mathcal{A}$：动作空间
    \item $P(s' | s, a)$：状态转移概率
    \item $R(s, a)$ 或 $R(s, a, s')$：奖励函数
    \item $\gamma \in [0, 1)$：折扣因子
\end{itemize}
\textbf{Markov 性}：下一个状态仅依赖当前状态和动作，与历史无关。
\end{importantbox}

\subsection{值函数与 Q 函数}

\begin{importantbox}{核心定义}
\textbf{状态值函数}：在状态 $s$ 下遵循策略 $\pi$ 的期望累积回报
$$V^\pi(s) = \mathbb{E}_\pi\left[\sum_{t=0}^{\infty} \gamma^t R(s_t, a_t) \mid s_0 = s\right]$$

\textbf{动作值函数}（Q 函数）：在状态 $s$ 下执行动作 $a$ 后遵循策略 $\pi$ 的期望累积回报
$$Q^\pi(s, a) = \mathbb{E}_\pi\left[\sum_{t=0}^{\infty} \gamma^t R(s_t, a_t) \mid s_0 = s, a_0 = a\right]$$

两者的关系：$V^\pi(s) = \sum_a \pi(a|s) Q^\pi(s, a)$
\end{importantbox}

\subsection{Bellman 方程}

\begin{importantbox}{Bellman 期望方程}
$$Q^\pi(s, a) = R(s, a) + \gamma \sum_{s'} P(s'|s,a) \sum_{a'} \pi(a'|s') Q^\pi(s', a')$$

\textbf{Bellman 最优方程}：
$$Q^*(s, a) = R(s, a) + \gamma \sum_{s'} P(s'|s,a) \max_{a'} Q^*(s', a')$$
\end{importantbox}

\subsection{策略迭代与价值迭代}

\begin{knowledgebox}{迭代视角的互补}
策略迭代（policy iteration）和价值迭代（value iteration）可以看作 Bellman 最优方程的两种实现路径：前者交替优化策略与价值函数，后者直接迭代最优价值。两者在收敛性质上都依赖于 Bellman 迭代的收缩性，但前者在每轮内对策略更新的精度更高，后者每步计算更简单。
\end{knowledgebox}

\begin{table}[H]
\centering
\begin{tabular}{ll}
\toprule
步骤 & 特征 \\
\midrule
策略迭代 & 明确构造行为策略（与旧策略做对齐）、值函数收敛后再改策略 \\
价值迭代 & 直接迭代 $V(s) \leftarrow \max_a [R+\gamma V(s')]$，省去显式策略存储 \\
两者的关系 & 价值迭代可看成策略迭代的一个特例（策略只在最后提取） \\
\bottomrule
\end{tabular}
\caption{在离散 MDP 中常见的迭代算子对比}
\label{tab:mdp-iteration}
\end{table}

\subsection{本章小结}

MDP 提供了 RL 的形式化框架。值函数和 Q 函数是评价策略好坏的核心工具，Bellman 方程描述了它们的递归关系。

